

\documentclass[12pt, letterpaper]{article}

% --- Paquetes básicos ---
\usepackage[spanish, es-tabla]{babel}
\usepackage[utf8]{inputenc}
\usepackage[T1]{fontenc}
\usepackage{newpxtext,newpxmath} % Fuente estilo Palatino
\usepackage[margin=3cm]{geometry}
\usepackage{titlesec}
\usepackage[autostyle]{csquotes}
\usepackage{amsmath}

% --- Para dibujar la V de Gowin ---
\usepackage{tikz}
\usetikzlibrary{positioning,babel,calc}

% --- Configuración de la sección (1. IDEA DE...) ---
\titleformat{\section}
  {\normalfont\large\scshape\centering} 
  {\thesection.}
  {0.5em}
  {}

\title
{
    Bitácora I \\[1ex] 
\large CA0305 - Herramientas de Ciencia de Datos II
}
\author{Anthonny Flores Rojas (C32975)\\ Andrey González Bastos (C33329)\\ Randal Picado Bermúdez (C36024)\\ Leonardo Vega Aragón (C38313) }
\date{\today}

%==============================
% CUERPO DEL DOCUMENTO
%==============================
\begin{document}
    
\maketitle


\newpage

\section{Tensor Broadcasting Basics}

En el contexto matemático, se define un tensor como un objeto algebraico
que describe la relación multilinear entre conjuntos de objetos algebraicos asociados
a un espacio vectorial. Esta sección afirma que en el contexto de \textit{Machine Learning} 
se acostumbra a manipular tensores con formas variadas. A su vez, presenta un mecanismo llamado \textit{Broadcasting}, que permite 
aplicar operaciones entre la i-ésima fila de un vector con la i-ésima fila de otro vector. Cabe mencionar, que dichos vectores podrían 
estar formados por otros vectores (como en las matrices, donde cada fila es un vector). Dicho mecanismo logra lo anterior al "estirar" el vector más pequeño (menor dimensión) para
que tenga las dimensiones del vector más grande.

El código asociado a esta sección permite sumar una constante a cada entrada de una matriz y hacer el producto de la í-ésima fila de un vector columna (el cual es un escalar) con la i-ésima fila
de la matriz (dicha fila es un vector). A grandes rasgos, se están haciendo varios productos escalares a varios vectores al mismo tiempo.  


\newpage


\section{Matrix Multiplication}

En esta sección se exploran dos formas de programar una función 
que calcule el producto de dos matrices (con las dimensiones adecuadas). La primera forma
consiste en usar únicamente las propiedades del paquete básico de Python, mientras que la otra 
forma utiliza únicamente las funciones optimizadas del paquete \textit{NumPy´s}. \\ 

Cabe destacar que este producto punto es el mismo que se estudió en el curso de Algebra Lineal, con todas
sus propiedades asociadas. \\ 

Para la primera función que se usaron los ciclos, se deben saber obviamente las dimensiones de las matrices para que se pueda hacer el producto matricial, 
luego se hace una matriz llena de ceros con esas dimensiones y se va llenando con el producto matricial entre las dos matrices tradicional que se estudió en Álgebra Lineal. \\ 

\begin{centering}

    C_{ij} = \sum_{k=1}^{n} A_{ik} \, B_{kj}

\end {centering} \\ 

Para la segunda función, se utilizó únicamente el \textit{np.dot()}, función de Numpy la cual permite obtener el produto matricial sin bucles y lo hace directamente. 

\newpage

\section{Element-wise Operations}

Las operaciones elemento a elemento son los componentes básicos de las redes neuronales. 
A diferencia de la multiplicación de matrices, que combina información de diferentes dimensiones, las operaciones elemento a elemento aplican la misma operación de forma 
independiente a cada par de elementos correspondiente. \\ 

Se definió una variable epsilon = $1e-8$ para poder hacer divisiones seguras. Se implementó la función que recibe las matrices A y B y hace las siguientes operaciones. \\ 

Suma: $a_{ij} + b_{ij}$, es decir suma entrada por entrada. \\ 

Multiplicación: $a_{ij} * b_{ij}$, es decir multiplica entrada por entrada. \\ 

División: $a_{ij} * b_{ij} + \epsilon$, es decir, divide entrada por entrada y usa el $\epsilon$ para evitar indeterminaciones. \\ 


\newpage

\section{Tensor Reshaping \& Transposing}

El redimensionamiento y la transposición de tensores son operaciones fundamentales en el aprendizaje profundo, especialmente cuando se trabaja con datos 
multidimensionales como imágenes, secuencias y lotes. Estas operaciones nos permiten: \\ 

\noindent Conversión de formato : Cambiar entre diferentes diseños de memoria (por ejemplo, NCHW frente a NHWC para imágenes). \\ \\ 
Manipulación de dimensiones : Reorganizar los datos para que sean compatibles con diferentes tipos de capas. \\ \\ 
Procesamiento por lotes : Reestructurar los datos aplanados para convertirlos nuevamente en lotes estructurados. \\ \\ 
Compatibilidad de API : Adaptar los formatos de entrada esperados para diferentes frameworks (PyTorch usa NCHW, TensorFlow usa NHWC). \\ \\  

La función que se implementó recibe un vector x de una sola dimensión. EL primer paso consiste en verificar si x cumple con los elementos necesarios para rellenar
el tensor. Si x pasa esa "prueba", ya se le hace el reshape con el B,C,H,W para hacerlo un tensor en 4D en formato NCHW y con el transpose finalmente se convierte en el formato NHWC
que era el objetivo. 

\newpage

\section{Reduction Operations}

Las operaciones de reducción se fundamentan en la idea de aplicar una función a conjuntos de elementos dentro de un tensor para obtener un valor representativo. 
Matemáticamente, esto implica tomar una colección de valores y aplicar operaciones como la suma, el promedio (suma entre 
la cantidad de elementos), el máximo y el índice del máximo.Estas operaciones reducen la dimensionalidad del tensor porque reemplazan múltiples valores por un resumen matemático,
permitiendo analizar y procesar datos de forma más eficiente.

\newpage

\section{Vector Norms}

Las normas vectoriales son medidas que permiten cuantificar la magnitud o tamaño de un vector en un espacio. La norma L1 se define 
como la suma de los valores absolutos de sus componentes, lo que implica medir la distancia considerando únicamente desplazamientos 
por cada eje. Por otro lado, la norma L2 se define como la raíz de la suma de los cuadrados de sus componentes, lo que corresponde a 
la distancia euclidiana desde el origen. Ambas normas reducen la información de un vector a un solo valor representativo, permitiendo 
comparar magnitudes y analizar datos.

\newpage

\section{Dot Product \& Cross Product}

El producto punto entre dos vectores se define como la suma del producto de sus componentes correspondientes, lo que permite medir 
el grado de alineación entre ellos. Matemáticamente, también se relaciona con el ángulo entre los vectores, siendo positivo cuando apuntan en direcciones 
similares, cero si son perpendiculares y negativo si apuntan en direcciones opuestas.

Por otro lado, el producto cruz se define como una operación que genera un vector perpendicular a los dos vectores originales. Su 
magnitud representa el área del paralelogramo formado por los vectores, y su dirección se determina mediante la regla de la mano 
derecha. Esta operación solo está definida en tres dimensiones y es clave en geometría y física.

\newpage

\section{Einstein Summation (einsum) Basics}

La notación de suma de Einstein (np.einsum) es una de las herramientas más potentes y expresivas de NumPy para operaciones con tensores. 
Proporciona una forma concisa y matemáticamente elegante de especificar contracciones, reducciones y transformaciones complejas de tensores mediante una notación de cadena sencilla.

Para esto, se implementó una función donde para la transpuesta, suma de todos los elementos de la matriz, suma de  filas, suma de columnas y producto matricial, se hace el juego de los 
índices. Por ejemplo en la transpuesta, basta con intercambiar los i con los j. Aprender este juego con los índices con la función einsum hace mucho más eficientes muchos procesos. 


\newpage

\section{Advanced Einsum (Batch MatMul)}

La multiplicación de matrices por lotes es fundamental para el aprendizaje profundo moderno, especialmente en las arquitecturas Transformer, donde los mecanismos 
de atención requieren operaciones por lotes eficientes. Este problema demuestra la einsumelegancia con la que se manejan las complejas contracciones 
de tensores multidimensionales, que serían engorrosas con las operaciones matriciales estándar. \\ 


Dados dos tensores $Q, K \in \mathbb{R}^{B \times H \times S \times D}$, se calcula el producto $QK^\top$ para cada par $(b, h)$:

\[
\text{Scores}_{bhij} = \sum_{d=0}^{D-1} Q_{bhid} \cdot K_{bhjd}
\]

El índice $d$ se suma (contrae); los índices $b, h, i, j$ se conservan en la salida $\text{Scores} \in \mathbb{R}^{B \times H \times S \times S}$.

En notación \texttt{einsum}:

\[
\texttt{"bhid, bhjd -> bhij"}
\]

Equivalente para un solo par $(b, h)$:

\[
\text{Scores}[b,h] = Q[b,h] \cdot K[b,h]^\top \quad (S \times D)\cdot(D \times S) = (S \times S)
\]

Se implementó la que recibe dos tensores Q y K de forma (B, H, S, D) y usa einsum con la cadena 'bhid,bhjd->bhij' para calcular el producto punto entre 
cada query y cada key. El índice d desaparece porque se suma, mientras que b, h, i, j se conservan, dando una salida de forma (B, H, S, S).

\newpage

\section{Gradient of Sum}

El gradiente por definicion es un vector 

\[
\nabla f(x) = \left( \frac{\partial f}{\partial x_1}, \frac{\partial f}{\partial x_2}, \dots, \frac{\partial f}{\partial x_n} \right)
\]
el cual se toma como generalización multiplicable de la derivada convencional, que mide la dirección en la cual se da el mayor crecimiento (o decrecimiento en su negativa) con un vector que es ortogonal a la superficie de la función (siendo esta función diferenciable).

Desde un punto de vista matemático, el algoritmo empieza de la suposición de que dado que las derivadas parciales de \(\displaystyle y = \sum_i x_i\) son iguales idénticas, ya que estas son funciones lineales; por regla de la cadena se tiene que 
\[
\frac{\partial L}{\partial x_i} = \frac{\partial L}{\partial y} \frac{\partial y}{\partial x_i}  = \frac{\partial L}{\partial y}
\]
en cuyo caso \(\frac{\partial L}{\partial x_i}  = \text{grad\_y}\). 

Gracias a que las derivadas parciales son iguales entre si, entonces el gradiente es simplemente el gradiente de y en cada una de las componentes del gradiente. Basta con tomar una matriz de dimensiones dadas y llenar todas las componentes con el valor del gradiente de y.


\newpage

\section{Gradient of Matrix Multiplication}



\newpage

\section{One-Hot Encoding}

\newpage

\section{Softmax Implementation (stable)}

En esta sección se destaca la relevancia de la función \textit{Softmax} en el contexto de redes neuronales. Cuando 
una red neuronal hace un análisis, devuelve un \textit{logit}, el cual se puede entender como un vector que guarda información acerca 
de la respuesta al problema que la red neuronal estaba analizando. Este vector por si solo no representa una distribución de probabilidad, pero 
gracias a la función \textit{Softmax} se puede generar un vector que guarde la información del \textit{logit} y que al mismo tiempo sea una distribución 
de probabilidad válida. Todo lo anterior permite interpretar los resultados del modelo como probabilidades, lo que permite cuantificar los errores del modelo.

Posteriormente se mencionan varios puntos que sirven para ejemplificar la importancia de está función en general, pero se pueden sintetizar en que \textit{Softmax} permite
conectar los resultados obtenidos por una red neuronal con predicciones interpretables. 

Matemáticamente, \textit{Softmax} es una matriz de dimensión $N\times C$ con las mismas dimensiones que la matriz original obtenida a partir del modelo, donde cada fila representa una muestra y cada columna representa una clasificación. Así, la matriz asociada a \textit{Softmax} está dada por:

$$\textit{Softmax}(x)_i = \frac{e^{x_i-\max(x)}}{\sum_{j=1}^C e^{x_j}} $$

Donde $x_i$ es el elemento ubicado en una cierta entrada de la matriz original, $max(x)$ es el máximo de la fila asociada al elemento mencionado anteriormente y $\sum_j=1^C e^x_j$ es la suma de las exponenciales asociadas a los elementos que están en la misma fila que el $x_i$ mencionado anteriormente. 

Primeramente, al aplicar la operación $\exp(x)$ a cada entrada de la matriz, entonces se asegura que todas las entradas de la nueva matriz van a ser positivas (lo cual es una condición necesaria para que un número represente una probabilidad). Por último, al definir una nueva matriz 
como la matriz anterior pero donde la iésima está multiplicada por el inverso de la suma de los elementos de dicha fila, gracias a esto es fácil comprobar que la suma de todos los elementos de la iésima fila de la nueva matriz será igual a 1, lo que es necesario al momento de definir una función de densidad (condición de normalización).
Además, es posible mostrar que: \begin{enumerate}
    \item ($Orden$) Dadas dos entradas de la matriz original $x_i$, $x_j$, si se tiene que $x_i > x_j$, entonces $\textit{Softmax}(x)_i>\textit{Softmax}(x)_j$ (pues la exponencial es creciente). 
    \item ($Diferenciabilidad$) Al ver la nueva matriz como una función $\textit{Softmax}(x)$, entonces se tiene que es infinitamente diferenciable entrada por entrada.
    \item ($Invarianza$) Sumar una constante a todas las entradas de una misma fila de la matriz original no cambia a la matriz \textit{Softmax} asociada.

Sin embargo, la función $\textit{Softmax}(x)$ tiene el problema de que las funciones exponenciales que definen sus entradas podrían causar un \textit{overflow} muy fácilmente (pues las exponenciales crecen súmamente rápido). Para solucionar esto, se puede aprovechar la propiedad de 
invarianza descrita anteriormente para acotar dichas exponenciales. Esto se puede lograr al restar el máximo elemento de una fila a cada elemento de la fila. De esta manera, se asegura que todos los elementos de la fila van a ser no positivos. Esto permite que, 
posteriormente, al aplicar la función $\exp(x)$ a cada entrada de la matriz, dicha función esté acotada 
superiormente (por lo que aplicar dicha función no va a requerir tanto poder computacional, pues no tendrá que 
calcular números más grandes que 1). 
\end{enumerate}





\newpage

\section{Cross-Entropy Loss (Manual)}

Se describe la función de pérdida \textit{Cross-entropy loss} como la función más usada al momento de 
realizar tareas de clasificación en el contexto de redes neuronales. Esta función permite cuantificar el error entre 
la función de distribución de probabilidad obtenida con el modelo y la distribución verdadera. Esto es fundamental, por ejemplo, al momento de 
identificar imágenes. Sin esta función, las redes neuronales no podrían ser entrenadas para clasificar datos de manera eficiente.

Esta función requiere dos parámetros: \begin{enumerate}
    \item \textit{Predicted distribution:} Un vector de la forma $p=[p_1, p_2, \dots, p_C]$ obtenid a partir de la función \textit{softmax}.
    \item \textit{True distribution:} Un vector tal que identifica la "característica verdadera" (vale 1.0 para la característica verdadera, y 0 para las demás).
\end{enumerate}

Así, para una característica $y$ se define la función \textit{cross-entropy} como: 
$$H(p,y)=-log(p_y)$$ El logarítmo aparece con el fin de medir la entropía (llamada informalmente como la medida de sorpresa de un resultado). Así, una probabilidad baja implica una "mayor sorpresa" (mayor pérdida).

Así, dadas $N$ muestras, se define la pérdida total como: 
$$\text{Loss } = -\frac{1}{N}\sum_{i=1}^N \log(p_{i,y_i}+\varepsilon)$$

donde $p_{i,y_i}$ es la probabilidad que el modelo le asigna a la característica correcta $y_i$ para la i-ésima muestra. Note que la pérdida total es simplemente el promedio de las pérdidas. 
Note que dentro del logaritmo hay una constante $\varepsilon>0$, esto con el fin de evitar que el logarítmo se evalúe en 0 (esto es necesario pues, si $p_{i,y_i}=0$ entonces $\log(p_{i,y_i})\rightarrow -\infty$, lo que causaría un \textit{overflow}).
 

\newpage

\section{Numerical Stability (Log-Sum-Exp)}

En esta sección se describe un método para calcular números de la forma: 
$$\log\left( \sum_{i=1}^n e^{x_i} \right)$$
con el fin de evitar un \textit{overflow}, pues las funciones exponenciales son reconocidas por crecer sumamente rápido.

Esto último es relevante pues, de acuerdo a la información brindada por la página \textit{Papercode}, la expresión anterior
suele aparecer con gran frecuencia en el contexto de \textit{Machine Learning}, por ejemplo al aplicar el logaritmo a una función \textit{softmax}. 

El método consiste en aprovechar la siguiente identidad matemática: 

$$\log\left( \sum_{i=1}^n e^{x_i} \right) = \max(x) + \log\left( \sum_{i=1}^n e^{x_i-\max(x)} \right)$$

Donde $\sum_{i=1}^n e^{x_i}$ es la suma de las función $exp(x)$ aplicada a los elemento de la i-esima fila de una matriz. Luego, $\max(x)$ es el máximo de los elementos de la fila mencionada anteriormente.

\textbf{Demostración identidad}: $\log\left( \sum e^{x_i} \right)=\log\left( \sum e^{x_i-\max(x)+\max(x)} \right)=\log\left( e^{\max(x)} \sum e^{x_i-\max(x)} \right)=\max(x) + \log\left( \sum_i=1^n e^{x_i-\max(x)} \right)$

Así, note que los elementos sobre los cuales se están aplicando la función $\exp(x)$ son no positivos, lo que implica que cada exponencial está acotada por 1 (lo que impide que haya un \textit{overflow}).


    
\end{document}
